\documentclass{article}
\usepackage{graphicx} % Required for inserting images
\usepackage{caption}
\usepackage[style=numeric-comp, sorting=none]{biblatex}
\usepackage[serbian]{babel}
\usepackage{hyperref}
\usepackage{xurl}
\usepackage{enumitem}

\title{Pregled tema \\[0.3cm] \large kurs: Računarstvo i društvo \\[0.3cm] \large profesor: Sana Stojanović Đurđević}
\author{Tea Pavlović 299/2021}
\date{Septembar 2026}
\renewcommand*\contentsname{Sadržaj}

\begin{document}

\maketitle

\newpage

\tableofcontents

\newpage

\section{Istorijski pregled razvoja računarstva}

\begin{itemize}
    \item \textbf{Rana istorija}
        \begin{itemize}[label=-]
                \item Računaljka abakus, paskaline, mehaničke mašine: mehanički razboj za tkanje predstavlja prvu programabilnu mašinu.
                \item Čarls Bebidž dizajnira prvu programabilnu računsku mašinu, a Ada Bajron piše prve programe za nju.
                \item Električni telegraf, telefon, radio.
            \end{itemize}
    \item \textbf{Bušene kartice}
        \begin{itemize}[label=-]
            \item Herman Holerit kreira mašinu za čitanje bušenih kartica koja se koristila za obradu rezultata popisa stanovništva.
            \item Od Holeritove kompanije kasnije nastaje čuvena kompanija IBM.
        \end{itemize}
    \item \textbf{Elektronski računari}
        \begin{itemize}[label=-]
            \item ABC je jedan od prvih elektronskih računara specijalne namene.
            \item Računar Kolos kreiran za razbijanje nemačke šifre zasnovane na mašini Enigma.
            \item ENIAC je bio prvi elektronski računar opšte namene.
            \item Fon Nojmanova arhitektura omogućava proizvodnju računara sa skladištenim prostorima.
            \item Definiše se jasna podela između hardvera i softvera, a osnovni elementi računara postaju procesor i glavna memorija.
        \end{itemize}
    \item \textbf{Pojava interneta}
        \begin{itemize}[label=-]
            \item ARPANET kreiran da se povežu vladini računari u SAD-u.
            \item TCP/IP protokoli nastaju kako bi se uključile i druge slične mreže.
        \end{itemize}
    \item \textbf{Presonalni računari}
        \begin{itemize}[label=-]
            \item Nastaje prvi Intel mikroprocesor koji oogućava proizvodnju malih, ali dovoljno brzih, kućnih računara.
            \item Prvi već sklopljeni računar bio je Apple.
        \end{itemize}
    \item \textbf{World Wide Web}
        \begin{itemize}[label=-]
            \item Berners Li kreira sistem za povezivanje informacija preko linkova, a potom nastaju i veb pregledači.
        \end{itemize}
\end{itemize}

Računarstvo se još uvek naglo razvija, a interesantne nove tehnologije obuhvataju mobilne uređaje, društvene mreže, blokčejn, veštačku inteligenciju, virtuelnu realnost itd.

\textbf{Predlog teme:} 

Quantum Computing: Uticaj kvantnih računara na svakodnevni život

\url{https://www.jstor.org/stable/pdf/48784776.pdf?refreqid=fastly-default%3A147b9b0732dd179f9607885641127b2f&ab_segments=&initiator=&acceptTC=1}

\url{https://ietresearch.onlinelibrary.wiley.com/doi/full/10.1049/iet-qtc.2020.0026}

\textbf{Primeri seminarskog rada iz repozitorijuma:}

Predrag Mitić - Informacione tehnologije kroz evoluciju društva (2022)

Mila Lukić - Ravnopravnost žena u računarstvu i informatici (2024)

\section{Etika}

Etika je grana filozofije koja se bavi izučavanjem morala, racionalnom sistematizacijom ljudskih verovanja i ponašanja, kao i odbranom i preporučivanjem koncepata ispravnog i pogrešnog ponašanja. Moral predstavlja pravila šta ljudi mogu, a šta ne mogu da rade u različitim situacijama. Podoblasti etike su metaetika, normativna etika i primenjena etika.

\textbf{Etičke teorije:}

\begin{enumerate}[label=\textbf{\arabic*.}]
    \item \textbf{Subjektivni relativizam} - Svaki pojedinac za sebe odlučuje šta je dobro, a šta loše.
    \item \textbf{Kulturološki relativizam} - Značenje dobrog i lošeg zavisi od principa koje definiše konkretno društvo, pa su ova pravila podložna promenama kroz vreme
    \item \textbf{Teorija duhovnog komandovanja} -Bazirana na ideji da je nešto dobro ako je u skladu sa Božijom voljom, a loše ako je u suprotnosti sa njom.
    \item \textbf{Etički egoizam} - Moralno dobrom akcijom se smatra ona koja pojedincu dugoročno donosi maksimalnu dobit.
    \item \textbf{Kantijanizam} - Moralna pravila koja su univerzalna za sve ljude.
    \item \textbf{Utilitarizam postupka} - Akcija je dobra ako je dobrobit akcije za sve uključene strane veća od štete.
    \item \textbf{Utilitarizam pravila} - Moralna pravila koja kada bi ih svi pratili povećavaju ukupnu sreću svih uključenih strana.
    \item \textbf{Teorija socijalnog ugovora} - Skup pravila, odnosno zakona, koji određuju kako se ljudi međusobno ponašaju.
    \item \textbf{Etika vrline} - Akcija je moralno ispravna ako bi je u istoj situaciji odabrala osoba koja poseduje vrline, to jest osobine koje su potrbne ljudskim bićima da bi napredovala i bila istinski srećna.
\end{enumerate}

Etički stavovi su  tehnicizam, optimizam i skepticizam. Ovi stavovi se odnose na etički pogled na tehnološki razvoj. 

Računarska etika je posebna grana etike koja se bavi etičkim pitanjima vezanim za informacione sisteme. Ipak, zvanična definicija ove oblasti još uvek ne postoji.

Pristupi računarskoj etici:

\begin{enumerate}[label=\arabic*.]
    \item Pristup nepostojanja - računarska etika ne treba da postoji.
    \item Profesionalni pristup - striktno profesionalna etika koja obuhvata probleme rada profesionalaca.
    \item Radikalni pristup - nova, jedinstvena grana etike koja zahteva nove pristupe i principe.
    \item Konzervativni pristup - računarska etika je podgrana primnjene etike gde uvodimo nove oblike starih problema.
    \item Inovativni pristup - podgrana primenjene etike, ali se uvode novi principi rešavanja problema.
\end{enumerate}

\textbf{Predlog teme:}

Trolley Problem: Etika autonomnih vozila

\url{https://pmc.ncbi.nlm.nih.gov/articles/PMC10846950/pdf/rsos.231393.pdf}

\url{https://arxiv.org/pdf/1510.03346}

\textbf{Primeri seminarskog rada iz repozitorijuma:}

Jovan Milić - Etika korišćenja veštačke inteligencije za prepoznavanje emocija u dečijim igračkama (2022)

Nemanja Badrić - Analiza sistema Gotham: Integracija velikih podataka i etika AI u vojnim i obaveštajnim sistemima (2026)

\section{Upotreba interneta, spam poruke}

Spam poruke predstavljaju neželjenu elektronsku poštu koja može u velikoj meri zatrpati sanduče korisnika. Sadržaji ovakvih poruka su najčešće reklame, ponude, eksplicitni sadržaj i prevare koje prevashodno čini krađa identiteta. Prednost spama je ta što je izuzetno jeftin, pa je čest izbor marketinških firmi. Na takve poruke ne treba odgovarati, pre svega zbog bezbednosti, a i kako pošiljaoc ne bi sa sigurnošću znao da je vaš mejl aktivan i time slao sve više ovakvog sadržaja. Spam može imati i negativan uticaj na kompanije i njihove zaposlene. Oni trpe posledice smanjenja propusne moći Interneta, zauzimanja resursa poput memorije i smanjenja produktivnosti zaposlenih koji troše vreme na filtriranje poruka.
Većina modernih imejl servisa ima dobre spam filtere koji prepoznaju spam poruke i direktno ih šalju u poseban folder kako bi nas upozorili da te poruke verovatno ne želimo da otvaramo. Nažalost, dešava se i da u ovom folderu povremeno završe poruke koje su nam važne ili koje mi šaljemo drugima.

Internet i veb nisu isto. Internet se sastoji od mrežne infrastrukture i internet protokola (IP), dok World Wide Web (WWW), ili skraćeno samo veb, predstavlja uslugu koja nam dostavlja međusobno povezane dokumente i oslanja se na HTTP protokol.

\textbf{Predlog teme:}

Digital Archives: Problemi očuvanja zastarelih podataka

\url{https://macsphere.mcmaster.ca/items/a1257ee3-e800-426a-8595-2624a1750e4e}

\url{https://www.cambridge.org/core/journals/german-law-journal/article/right-to-be-forgotten-in-the-digital-age-the-challenges-of-data-protection-beyond-borders/3E3E182352F1AD555CBB788E2380E23F}

\url{https://d1wqtxts1xzle7.cloudfront.net/52257569/10579_2004_Article_153071-libre.pdf?1490201714=&response-content-disposition=inline%3B+filename%3DDigital_preservation_a_time_bomb_for_dig.pdf&Expires=1788815074&Signature=HrnsuQeUW4wZCwfHKxxznQNc~7enKUTLQPbxZ60XMNDFeeG7ecyoHdF58AYLv8x~qVJz5uvDvFEo9t4kuZQeFiNE3QhAkFYGlM0N7FPCQlyy4-plJ0HHUlr8HeEwawuIGoRaTmL2lut~W67X5g5GtiiK1CwA5p9y0JZvqZulz3u~ypMwujRpMQTeWnI2o3kTHLJg8s~FQiLuZC7Higk-Zln0BnaITRJMMz1r-amOWKwqeuTUJUJnCR~~ZKliGqIu65uWkmmF2cJXsLexcpKMOnr00ym0vrVEAaLNiDgOuSHuPUnr0h42dLjV1FOXyb-jAS0iMK~1Lg41li4WBCv8vg__&Key-Pair-Id=APKAJLOHF5GGSLRBV4ZA}

\textbf{Primeri seminarskog rada iz repozitorijuma:}

Miodrag Todorović - Internet i nove tehnologije kao pomoć u rešavanju problema anksioznosti i emotivnih okidača (2023)

Luka Bura - Should knowledge be free (2023)

\section{Cenzura i sloboda govora, deca i neprikladan sadržaj}

Cenzura je pokušaj da se reguliše ili spreči javni pristup informacijama ili materijalima koji se smatraju štetnim, opasnim ili uvredljivim.

\textbf{Cenzura kroz istoriju:}

\begin{itemize}
    \item Rimska vlada - koreni cenzure
    \item Inkvizicija u srednjem veku - najveći pokret cenzure od strane vlasti i crkve
    \item Pojava štamparske mašine - otežava cenzuru
\end{itemize}

Cenzuru možemo podeliti na direktnu i samocenzuru.

\textbf{Primeri:}

\begin{itemize}
    \item Monopol nad informacijama od strane vladinih institucija - Sovjetski savez, Kina, Saudijska Arabija, Srbija uprkos zabrane ustavom.
    \item Pregled publikacije - većina država zabranjuje ili barem proverava sadržaj vezan za osetljiva pitanja.
    \item Licenciranje i registracija - radio i televizijske stanice.
\end{itemize}

Samocenzura je sveprisutna na društvenim mrežama jer sami odlučujemo koji sadržaj objavljujemo.

Sloboda govora je u direktnoj suprotnosti sa cenzurom, pa u mnogim slučajevima cenzuru smatramo negativnom. Ipak, deljenje sadržaja koji šire mržnju, maliciozne savete ili dezinformacije može dovesti do užasnih posledica po one kojima je takav sadržaj upućen. Ovakvi sadržaji se često nazivaju i neprikladni, a posebno mogu naškoditi deci koja sve mlađa dobijaju pristup internetu. Primeri: pornografski i nasilni sadržaji, sajtovi za klađenje itd. Osim roditeljske kontrole i razgovora sa decom, postoji softver poput veb filtera koji sprečavaju određenim veb stranicama da se prikažu u pregledaču.

\textbf{Predlog teme:}

Shadowbanning: Algoritamska cenzura

\url{https://onlinelibrary.wiley.com/doi/full/10.1002/poi3.287}

\url{https://dl.acm.org/doi/abs/10.1145/3701191}

\textbf{Primeri seminarskog rada iz repozitorijuma:}

Danilo Nikolaš - Kako video igre utiču na malu decu (2025)

Mila Gligorić - Sistemski gubitak podataka u velikim jezičkim modelima: kolaps modela i algoritamska diskriminacija (2026)

\section{Internet prevare i zavisnost od interneta}

Internet prevare obično podrazumevaju krađu identiteta, pecanje i Nigerijsku prevaru. 

Krađa identiteta je zloupotreba identiteta druge osobe, odnosno njenog imena, ličnih dokumenata, matičnog broja, broja bankovnog računa i kartica. 

Pecanje je najlakši način da se dođe do identiteta ili novca žrtve. Lopov najčešće šalje poruku koja izgleda kao da ju je poslala neka poznata kompanija ili osoba, sa ciljem da se žrtva „upeca“ tako što će kliknuti na hiperlink i upisati tražene informacije koje potom direktno stižu do lopova. Najčešći mediji za pecanje su imejl i društvene mreže. Ovaj vid prevare je lakše izvesti ako žrtva internetu uglavnom pristupa preko mobilnog uređaja jer su ekrani manji i teže je primetiti anomaliju u imejl adresi pošiljaoca ili stranici kojoj pristupamo. 

Nigerijska prevara je nastala još 1980. godine i podrazumeva situaciju kada se od žrtve traži pomoć za transfer novca.

Zavisnost od interneta je sve češća i predstavlja sve veći problem. Teško ju je prepoznati jer se dosta razlikuje od drugih zavisnosti, a neki znaci mogu biti opravdani, na primer, zbog posla. Prvi test zavisnosti od interneta je uveden 2005. godine i danas se smatra relevantnim. 

Posledice ove zavisnosti uključuju i mentalne i fizičke probleme: anksioznost i depresiju, gubljenje percepcije o vremenu, glavobolje, slabljenje vida itd. Kod mladih osoba ove i druge posledice su sve izraženije, a najčešći vid zavisnosti od interneta je zavisnost od društvenih mreža koja je klasifikovana kao poremećaj samo u Kini i Južnoj Koreji. Zavisnost od igrica je takođe noviji oblik zavisnosti, a javlja se najčešće pri igranju igrica u kojima se kreiraju virtuelni avatari i igrači zajedno borave u virtuelnom svetu. Kao i klasično klađenje, online klađenje takođe izaziva zavisnost, a pritom je mnogo dostupnije.

\textbf{Predlog teme:}

Dark Patterns in Social Media: Dizajn koji razvija zavisnost

\url{https://www.frontiersin.org/journals/public-health/articles/10.3389/fpubh.2021.641673/full?trk=public_post_comment-text}

\url{https://dl.acm.org/doi/abs/10.1145/3563657.3595964}

\textbf{Primeri seminarskog rada iz repozitorijuma:}

Luka Miletić - Zavisnost od pametnih telefona i njihov uticaj (2022)

Sofija Janevska - Algoritamska manipulacija pažnjom: tehnike i posledice maksimizacije korisničkog angažovanja (2026)

\section{Bezbednost interneta}

Haker je osoba koja neovlašćeno pristupa računarima i računarskim mrežama. Kako bismo se zaštitili treba koristiti jake šifre na nalozima.

Firesheep je ekstenzija za Firefox koja pomoću kolačića hakuje veb stranice koje nisu enkriptovane. Programer koji je napravio ovu ekstenziju imao je za cilj da skrene pažnju velikim kompanijama društvenih mreža koliko su korisnički profili bili nebezbedni u tom trenutku.

Zlonamerni softver je program dizajniran da nanese štetu ciljnom korisniku. Virusi su vrsta zlonamernog softvera koji prave štetu u vidu usporavanja procesora i drugih komponenti računara ili napadajući korisničke podatke i pokušavajući da preuzmu kontrolu nad računarom, a ime su dobili po načinu na koji se šire. Virus je deo drugog programa koji zovemo domaćin i sposoban je da pravi sopstvene kopije. Prenose se preko prenosivih diskova, CD-ova, skidanja podataka preko interneta i imejla.

Još neki oblici zlonamernog softvera su: računarski crvi, ukrštanje veb lokacija, usputno preuzimanje programa, trojanski konj, softver za špijuniranje, mreža botova.

Antivirus je softver kojim se štitimo od virusa i sprečavamo njihovo širenje. Ipak, neophodno ih je redovno ažurirati kako bi imali mogućnost da nas odbrane i od najnovijih virusa. Zaštitni zid je softver za regulisanje mrežnog saobraćaja i pomaže u odbrani od malicioznog softvera tako što blokira sumnjivi saobraćaj.

\textbf{Predlog teme:} 

Adversarial Machine Learning: Napadi na autonomne sisteme i pametna vozila

\url{https://dl.acm.org/doi/pdf/10.1145/2046684.2046692}

\url{https://dl.acm.org/doi/pdf/10.1145/3453158}

\textbf{Primeri seminarskog rada iz repozitorijuma:}

Aleksa Jovanović - Bezbednost naših lozinki (2024)

Katarina Grbović - Ofanzivna bezbednost i pentesting (2025)

\section{Privatnost na internetu}

Privatnost na internetu možemo zaštititi tako što ograničimo podatke koje postavljamo na društvene mreže, koristimo pregledače koji ne čuvaju kolačiće i istoriju, inkognito mod, VPN i kvalitetan antivirus. 

Kolačići, ili preciznije HTTP kolačići, su tekstualni dokumenti koji sadrže informacije o veb sajtovima. Oni imaju jedinstvene identifikatore i čuvaju informacije poput podataka o prijavi korisnika, personalizovanih podešavanja na sajtu ili sadržaja korpe. 

\textbf{Vrste kolačića:}

\begin{enumerate}[label=\arabic*.]
    \item Sesijski (privremeni) - ne čuvaju se nakon što korisnik zatvori veb sajt ili pregledač.
    \item Trajni kolačići prve strane - ostaju sačuvani, a koristi ih isključivo veb sajt koji ih je napravio.
    \item Trajni kolačići treće strane (kolačići za praćenje) - ostaju sačuvani, a koriste se za praćenje korisnika i prikupljanje informacija kroz korisničke aktivnosti na internetu.
\end{enumerate}

U prošlosti, kolačići su bili podrazumevano uključeni, a 2018. godine je uveden zakon koji obavezuje sajtove da informišu korisnike o svim kolačićima i dopuste im da biraju koje kolačiće žele da omoguće, osim onih neophodnih za funkcionisanje sajta. 

EU opšta uredba o zaštiti podataka (GDPR) je takođe stvorena radi kreiranja zakona o privatnosti podataka u Evropi i odnosi se na sva tela koja prikupljaju lične podatke.

\textbf{Predlog teme:}

Big brain data: Privatnost moždanih podataka

\url{https://link.springer.com/article/10.1007/s12152-018-9371-x}

\url{https://link.springer.com/chapter/10.1007/978-3-030-32692-0_16}

\textbf{Primeri seminarskog rada iz repozitorijuma:}

Anđela Niketić - Privatni nalozi na društvenim mrežama: Neophodnost ili ne? (2024)

Mitar Avramović - OSINT- Lupa u moru (2024)

\end{document}
